\section{Step 7: Tor / SOCKS5 Support}

Direct connections reveal your IP to the destination. With \texttt{-T}, we
instead open a TCP connection to a local Tor SOCKS5 proxy (default
\texttt{127.0.0.1:9050}) and ask \emph{it} to reach the real host. The HTTP
or HTTPS session then rides on that tunnel.

\subsection{Why SOCKS5 (not HTTP CONNECT)?}

Tor's default application proxy speaks SOCKS5. Domain names can be sent to
the proxy (\texttt{ATYP = 0x03}), so DNS resolution also happens inside Tor
when configured that way---important for anonymity.

\subsection{Minimal Handshake}

\begin{enumerate}[leftmargin=1.4em]
    \item \textbf{Greeting:} version~5, one method ``no authentication''
          (\texttt{0x00}).
    \item \textbf{Server reply:} version~5, method \texttt{0x00} accepted.
    \item \textbf{Connect request:} \texttt{CMD=CONNECT}, address type domain,
          hostname length + bytes, port in network byte order.
    \item \textbf{Reply:} \texttt{REP=0x00} means success. (We ignore the
          bind address fields for a minimal client.)
\end{enumerate}

\subsection{New Code --- \texttt{socks5\_connect}}

\begin{lstlisting}[style=snippet,language=C]
static int socks5_connect(int sock, const char *host, uint16_t port)
{
    unsigned char buf[512];
    size_t hlen = strlen(host);
    if (hlen > 255) return -1;

    /* greeting: VER=5, NMETHODS=1, METHOD=0 (no auth) */
    buf[0] = 0x05; buf[1] = 0x01; buf[2] = 0x00;
    if (write(sock, buf, 3) != 3) return -1;

    if (read(sock, buf, 2) != 2 || buf[0] != 0x05 || buf[1] != 0x00)
        return -1;

    /* CONNECT request with domain name */
    buf[0] = 0x05; buf[1] = 0x01; buf[2] = 0x00; buf[3] = 0x03;
    buf[4] = (unsigned char)hlen;
    memcpy(buf + 5, host, hlen);
    buf[5 + hlen] = port >> 8;
    buf[6 + hlen] = port & 0xff;

    if (write(sock, buf, 7 + hlen) != (ssize_t)(7 + hlen)) return -1;
    if (read(sock, buf, 4) != 4 || buf[1] != 0x00) return -1;
    return 0;
}
\end{lstlisting}

\subsection{Wiring Into \texttt{main}}

When Tor is enabled, DNS/\texttt{connect} targets the \emph{proxy}, not the
origin. After the TCP connect succeeds, run the SOCKS5 handshake with the
real host and port from \texttt{url\_t}. TLS (if any) still uses the real
hostname for SNI.

\begin{lstlisting}[style=snippet,language=C]
const char *connect_host = use_tor ? tor_proxy : u.host;
uint16_t connect_port = use_tor ? 9050 : u.port;

if (use_tor) {
    char pbuf[256];
    strncpy(pbuf, tor_proxy, sizeof(pbuf) - 1);
    char *colon = strrchr(pbuf, ':');
    if (colon) {
        *colon = '\0';
        connect_port = (uint16_t)strtol(colon + 1, NULL, 10);
        connect_host = pbuf;
    }
}
/* getaddrinfo(connect_host) + connect ... */
if (use_tor && socks5_connect(s, u.host, u.port) != 0) {
    close(s);
    return EXIT_FAILURE;
}
\end{lstlisting}

\FullSourceStart{urlgetstep07}
\begin{lstlisting}[style=oldcode,language=C,firstnumber=1]
/* urlget.c -- Step 7: Tor / SOCKS5 */
#include <stdio.h>
#include <stdlib.h>
#include <string.h>
#include <stdbool.h>
#include <stdint.h>
#include <unistd.h>
#include <sys/types.h>
#include <sys/socket.h>
#include <netdb.h>
#include <openssl/ssl.h>
#include <openssl/err.h>

typedef struct {
    bool https;
    char host[256];
    uint16_t port;
    char path[512];
} url_t;

static inline int parse_url(const char *arg, url_t *u)
{
    char buf[1024];
    if (strlen(arg) >= sizeof(buf)) return -1;
    strcpy(buf, arg);
    char *p = buf;
    u->https = false;
    if (strncmp(p, "https://", 8) == 0) { u->https = true; p += 8; }
    else if (strncmp(p, "http://", 7) == 0) { p += 7; }
    char *sep = strpbrk(p, "/?");
    char *hostport = p;
    if (sep) {
        if (*sep == '?') {
            size_t ql = strlen(sep);
            if (ql + 2 > sizeof(u->path)) return -1;
            u->path[0] = '/';
            memcpy(u->path + 1, sep, ql + 1);
        } else {
            if (strlen(sep) >= sizeof(u->path)) return -1;
            strcpy(u->path, sep);
        }
        *sep = '\0';
    } else strcpy(u->path, "/");
    char *host = u->host;
    size_t hmax = sizeof(u->host) - 1;
    uint16_t def = u->https ? 443 : 80;
    u->port = def;
    if (*hostport == '[') {
        char *cl = strchr(hostport, ']');
        if (!cl) return -1;
        size_t hl = cl - (hostport + 1);
        if (hl >= hmax) return -1;
        memcpy(host, hostport + 1, hl);
        host[hl] = '\0';
        if (*(cl + 1) == ':')
            u->port = (uint16_t)strtol(cl + 2, NULL, 10);
    } else {
        char *col = strrchr(hostport, ':');
        if (col) {
            size_t hl = col - hostport;
            if (hl >= hmax) return -1;
            memcpy(host, hostport, hl);
            host[hl] = '\0';
            u->port = (uint16_t)strtol(col + 1, NULL, 10);
        } else {
            if (strlen(hostport) >= hmax) return -1;
            strcpy(host, hostport);
        }
    }
    if (u->port == 0) u->port = def;
    return 0;
}
\end{lstlisting}
\begin{lstlisting}[style=newcode,language=C,firstnumber=last]
static int socks5_connect(int sock, const char *host, uint16_t port)
{
    unsigned char buf[512];
    size_t hlen = strlen(host);
    if (hlen > 255) return -1;

    buf[0] = 0x05; buf[1] = 0x01; buf[2] = 0x00;
    if (write(sock, buf, 3) != 3) return -1;
    if (read(sock, buf, 2) != 2 || buf[0] != 0x05 || buf[1] != 0x00)
        return -1;

    buf[0] = 0x05; buf[1] = 0x01; buf[2] = 0x00; buf[3] = 0x03;
    buf[4] = (unsigned char)hlen;
    memcpy(buf + 5, host, hlen);
    buf[5 + hlen] = port >> 8;
    buf[6 + hlen] = port & 0xff;
    if (write(sock, buf, 7 + hlen) != (ssize_t)(7 + hlen)) return -1;
    if (read(sock, buf, 4) != 4 || buf[1] != 0x00) return -1;
    return 0;
}
\end{lstlisting}
\begin{lstlisting}[style=oldcode,language=C,firstnumber=last]
static inline bool https_setup(int sock, const char *hostname,
                               SSL_CTX **out_ctx, SSL **out_ssl)
{
    SSL_CTX *ctx = SSL_CTX_new(TLS_client_method());
    if (!ctx) { ERR_print_errors_fp(stderr); return false; }
    SSL_CTX_set_verify(ctx, SSL_VERIFY_NONE, NULL);
    SSL *ssl = SSL_new(ctx);
    if (!ssl) { SSL_CTX_free(ctx); return false; }
    SSL_set_fd(ssl, sock);
    SSL_set_tlsext_host_name(ssl, hostname);
    if (SSL_connect(ssl) <= 0) {
        ERR_print_errors_fp(stderr);
        SSL_free(ssl); SSL_CTX_free(ctx);
        return false;
    }
    *out_ctx = ctx;
    *out_ssl = ssl;
    return true;
}

static inline void usage(const char *progname)
{
    fprintf(stderr,
        "Usage: %s [options] <URL>\n\n"
        "Options:\n"
        "  -A, --user-agent STR     Set User-Agent\n"
        "  -T, --tor [proxy]        Use Tor (default 127.0.0.1:9050)\n"
        "  -H, --header STR         Add custom header (repeatable)\n"
        "  -d, --data DATA          POST data (@file supported)\n"
        "  -o, --output FILE        Write output to file\n"
        "  -u, --user USER:PASS     Basic authentication\n"
        "  -I, --head               HEAD request\n"
        "      --no-location        Do not follow redirects\n"
        "      --max-redirects N    Max redirects (default: 10)\n",
        progname);
    exit(EXIT_FAILURE);
}

int main(int argc, char *argv[])
{
    const char *user_agent =
        "Mozilla/5.0 (X11; Linux x86_64; rv:128.0) "
        "Gecko/20100101 Firefox/128.0";
    const char *default_tor = "127.0.0.1:9050";
    const char *url_arg = NULL;
    const char *tor_proxy = NULL;
    bool use_tor = false;
    bool follow_redirects = true;
    int max_redirects = 10;
    bool head_request = false;
    char extra_headers[4096] = {0};
    char post_data[8192] = {0};
    char output_file[256] = {0};
    char auth_header[256] = {0};

    for (int i = 1; i < argc; i++) {
        if (strcmp(argv[i], "-A") == 0 ||
            strcmp(argv[i], "--user-agent") == 0) {
            if (i + 1 < argc) user_agent = argv[++i];
            else usage(argv[0]);
        }
        else if (strcmp(argv[i], "-T") == 0 ||
                 strcmp(argv[i], "--tor") == 0) {
            use_tor = true;
            if (i + 1 < argc && argv[i + 1][0] != '-' &&
                strchr(argv[i + 1], ':')) {
                tor_proxy = argv[i + 1];
                i++;
            }
        }
        else if (strcmp(argv[i], "-H") == 0 ||
                 strcmp(argv[i], "--header") == 0) {
            if (i + 1 < argc) {
                strncat(extra_headers, argv[++i],
                        sizeof(extra_headers) -
                        strlen(extra_headers) - 3);
                strcat(extra_headers, "\r\n");
            } else usage(argv[0]);
        }
        else if (strcmp(argv[i], "-d") == 0 ||
                 strcmp(argv[i], "--data") == 0) {
            if (i + 1 < argc) {
                const char *val = argv[++i];
                if (val[0] == '@') {
                    FILE *f = fopen(val + 1, "r");
                    if (f) {
                        fread(post_data, 1,
                              sizeof(post_data) - 1, f);
                        fclose(f);
                    }
                } else
                    strncpy(post_data, val, sizeof(post_data) - 1);
            } else usage(argv[0]);
        }
        else if (strcmp(argv[i], "-o") == 0 ||
                 strcmp(argv[i], "--output") == 0) {
            if (i + 1 < argc)
                strncpy(output_file, argv[++i],
                        sizeof(output_file) - 1);
            else usage(argv[0]);
        }
        else if (strcmp(argv[i], "-u") == 0 ||
                 strcmp(argv[i], "--user") == 0) {
            if (i + 1 < argc)
                snprintf(auth_header, sizeof(auth_header),
                    "Authorization: Basic %s\r\n", argv[++i]);
            else usage(argv[0]);
        }
        else if (strcmp(argv[i], "-I") == 0 ||
                 strcmp(argv[i], "--head") == 0)
            head_request = true;
        else if (strcmp(argv[i], "--no-location") == 0)
            follow_redirects = false;
        else if (strcmp(argv[i], "--max-redirects") == 0 &&
                 i + 1 < argc)
            max_redirects = atoi(argv[++i]);
        else if (url_arg == NULL)
            url_arg = argv[i];
    }

    if (url_arg == NULL) usage(argv[0]);
    if (use_tor && tor_proxy == NULL) tor_proxy = default_tor;

    url_t u = {0};
    if (parse_url(url_arg, &u) != 0) {
        fprintf(stderr, "Failed to parse URL\n");
        return EXIT_FAILURE;
    }
\end{lstlisting}
\begin{lstlisting}[style=newcode,language=C,firstnumber=last]
    const char *connect_host = use_tor ? tor_proxy : u.host;
    uint16_t connect_port = use_tor ? 9050 : u.port;
    char pbuf[256];
    if (use_tor) {
        strncpy(pbuf, tor_proxy, sizeof(pbuf) - 1);
        char *colon = strrchr(pbuf, ':');
        if (colon) {
            *colon = '\0';
            connect_port = (uint16_t)strtol(colon + 1, NULL, 10);
            connect_host = pbuf;
        }
    }
\end{lstlisting}
\begin{lstlisting}[style=oldcode,language=C,firstnumber=last]
    struct addrinfo hints = {0};
    hints.ai_family   = AF_UNSPEC;
    hints.ai_socktype = SOCK_STREAM;
    char ps[6];
\end{lstlisting}
\begin{lstlisting}[style=newcode,language=C,firstnumber=last]
    snprintf(ps, sizeof(ps), "%u", connect_port);

    struct addrinfo *res = NULL;
    if (getaddrinfo(connect_host, ps, &hints, &res) != 0) {
\end{lstlisting}
\begin{lstlisting}[style=oldcode,language=C,firstnumber=last]
        perror("getaddrinfo");
        return EXIT_FAILURE;
    }

    int s = -1;
    for (struct addrinfo *a = res; a; a = a->ai_next) {
        s = socket(a->ai_family, a->ai_socktype, a->ai_protocol);
        if (s < 0) continue;
        if (connect(s, a->ai_addr, a->ai_addrlen) == 0) break;
        close(s); s = -1;
    }
    freeaddrinfo(res);
    if (s < 0) { perror("connect"); return EXIT_FAILURE; }
\end{lstlisting}
\begin{lstlisting}[style=newcode,language=C,firstnumber=last]
    if (use_tor && socks5_connect(s, u.host, u.port) != 0) {
        close(s);
        return EXIT_FAILURE;
    }
\end{lstlisting}
\begin{lstlisting}[style=oldcode,language=C,firstnumber=last]
    SSL_CTX *ctx = NULL;
    SSL *ssl = NULL;
    if (u.https && !https_setup(s, u.host, &ctx, &ssl)) {
        close(s);
        return EXIT_FAILURE;
    }

    char req[4096];
    snprintf(req, sizeof(req),
        "GET %s HTTP/1.1\r\n"
        "Host: %s\r\n"
        "User-Agent: %s\r\n"
        "Accept: */*\r\n"
        "Connection: close\r\n"
        "\r\n",
        u.path, u.host, user_agent);

    if (u.https)
        SSL_write(ssl, req, strlen(req));
    else
        write(s, req, strlen(req));

    char buf[8192];
    ssize_t n;
    while ((n = u.https
                ? SSL_read(ssl, buf, sizeof(buf))
                : read(s, buf, sizeof(buf))) > 0)
        fwrite(buf, 1, n, stdout);

    if (u.https) {
        if (ssl) { SSL_shutdown(ssl); SSL_free(ssl); }
        if (ctx) SSL_CTX_free(ctx);
    }
    close(s);
    (void)follow_redirects; (void)max_redirects;
    (void)head_request; (void)extra_headers;
    (void)post_data; (void)output_file; (void)auth_header;
    return 0;
}
\end{lstlisting}
\FullSourceStop

Requires a running Tor daemon (or any SOCKS5 proxy) on the configured
address:

\begin{lstlisting}[style=snippet,language=bash]
./urlget -T https://check.torproject.org/
./urlget -T 127.0.0.1:9050 https://example.com/
\end{lstlisting}
